\documentclass{ULBreport}
\usepackage[utf8]{inputenc}
\usepackage[T1]{fontenc}
\usepackage{geometry}
\usepackage{graphicx}
\usepackage{float}
\usepackage{hyperref}
\usepackage[english,indonesian]{babel}

\title{Optimasi Pemrosesan Template LaTeX Menggunakan Pendekatan Berbasis AI}
\author{Academic Support Agent}
\date{\today}

\begin{document}
\maketitle

\begin{abstract}

\end{abstract}

\chapter{PENDAHULUAN}


\chapter{TINJAUAN PUSTAKA}


% Original content (commented out):
% \documentclass{ULBreport}

\sceau{Pictures/sceauULB.jpg}

\addbibresource{biblio.bib}

\input{common/packages}

\makeglossaries

% \newglossaryentry{strip mining}
% {
%   name=strip mining,
%   description={is the practice of mining a seam of mineral, by first removing a long strip of overlying soil and rock}
% }



% \newacronym{t}{t}{tons}



\begin{document}

\titleULB{
	title={Two bars and one elastic},
	course={Control System Design},
	author={Roberto ACQUAVIVA \\ other names},
	date={2021},
	teacher={Prof. Emanuele Garone},
	logo={Pictures/logos.jpg}
}

\listoftables % ToC for tables

\listoffigures % ToC for figures

\input{content/chapters.tex}

\printglossary

\printglossary[type=\acronymtype]

%Bibliography
\nocite{*}
\printbibliography[type=article,title=Articles]

\end{document}


\end{document}